\section{Design Principles} \label{sec:design-principles}

This chapter presents the derivation and evaluation of the design principles that constitute the primary contribution of this thesis. Section~\ref{ssec:derivation-methods} details the methodology used to derive the design principles from the results of the empirical study described in Chapter~\ref{sec:artefact-design}, and the evaluation and revision of the resulting principles through expert interviews. The coding results and first iteration of the design principles are presented in Section~\ref{ssec:initial-design-principles}. The results of the expert interviews and the resulting revisions to the design principles are reported in Section~\ref{ssec:design-principle-revision}. Finally, Section~\ref{ssec:uncovered-design-space} addresses a portion of the design space that could not be covered empirically, discusses potential reasons for this gap, and suggests theory-based design principles to address the gap in the design space.

\subsection{Derivation Methodology and Expert Interviews} \label{ssec:derivation-methods}

The design principles presented in this chapter are constructed following the grounded-theory approach \parencite{Gioia2013}. The grounded theory approach was chosen because its progressive abstraction from first-order concepts to second-order themes and aggregate dimensions is uniquely suited to the derivation of design principles from empirical data. Design principles are suited to provide the final abductive step in the grounded theory process \parencite{Gioia2013}, by connecting the inductively derived data structure to the pre-existing theoretical framework of \acp{MR} and \acp{DR} established in Section~\ref{ssec:meta-requirements}.

Coding was performed in two cycles based on the guidance presented in \textcite{Saldana2015}. The first cycle applied process codes \parencite{Corbin2014}, as the analysis focused on the interaction between participants and the \ac{AI} system. Decomposing this interaction into a sequence of discrete actions allowed for the identification of patterns in the participants' behaviour and their engagement with the system. In addition, each process code was matched with an in vivo code \parencite{Corbin2014} from the transcripts to preserve a direct connection to the underlying data for the subsequent abstraction levels. Following the first cycle coding, the raw codes were cleaned and consolidated by correcting typographical and minor wording variations without altering the underlying coding decisions. The second cycle coding applied pattern coding \parencite{Miles2013}, combining first-cycle codes of similar or related meaning into smaller sets of pattern codes to reduce the overall code space to a size suitable for further thematic analysis.

The pattern codes produced in the second cycle were used as first-order concepts, which were then grouped into second-order themes and aggregate dimensions. The second-order themes and aggregate dimensions were then used to derive the initial design principles using the \acp{MR} and \acp{DR} defined in Section~\ref{ssec:meta-requirements} as guidance. This inductive process of stepwise abstraction allows for the constructed data structure to be grounded in the empirical data. The use of the \acp{MR} and \acp{DR} as a guiding framework places the resulting design principles as connections between the inductively derived data structure and the pre-existing theoretical framework \parencite{Gioia2013}. Using the same codes between the two groups, also allows for a limited degree of quantitative comparison between the groups, complementing the qualitative construction of design principles.

The resulting design theory was then evaluated through expert interviews with \ac{HCI} and \ac{HAII} researchers and industry practitioners. The interviews were conducted as semi-structured interviews, aimed to evaluate each design principle individually as well as the design theory as a whole. Interview questions were designed to evaluate the design principles against the criteria of \textbf{accessibility}, \textbf{importance}, \textbf{novelty and insightfulness}, \textbf{actability and guidance}, and \textbf{effectiveness} proposed by \textcite{Iivari2021}. The design theory as a whole was evaluated for \textbf{deficits}, \textbf{redundancy}, \textbf{overload}, and \textbf{excess} \parencite{Janiesch2020}. The feedback gathered during the interviews was used to refine and revise the design principles.

\subsection{Initial Design Principles} \label{ssec:initial-design-principles}

% TODO: add citation for evaluation repo
The original transcripts from the study contained only speaker annotations and stage directions. To better support the construction of the grounded theory data structure, they were reformatted into a semi-structured format. It assigns each statement a unique identifier, used throughout this section to reference specific statements in the transcripts. These identifiers follow the format \texttt{[participant ID].[turn ID].[statement ID]}. Details on the format, including the additional metadata it carries, are documented in the supplementary materials \parencite{PLACEHOLDER}. The reformatted transcripts were used to construct the Gioia data structure through the following steps:

\begin{enumerate}
    \item \textbf{First-cycle Coding:} Transcripts were coded in sequence in a single pass of process coding and in vivo coding, resulting in 1250 first-cycle codes across 1466 coded statements. Each statement received one or more process codes, with the accompanying in vivo code serving as its empirical grounding.

    \item \textbf{Code Cleaning:} The first-cycle codes were deduplicated by unifying spelling and capitalisation differences. Additionally, codes that were in a direct sub- or supercode relationship were merged, resulting in 1163 remaining canonical codes.

    \item \textbf{Second-cycle Coding:} The pattern coding combined codes based on the primary verbs of the process codes. Initially, codes were grouped by their primary verbs and merged within each group. For example, codes like \enquote{Restating task goal} and \enquote{Restating problem goal} were merged on a code-by-code basis, because they describe the same action in the context of the experiment. When the verb groups reached a size of eight codes or fewer, the remaining codes were merged by grouping in semantic blocks\footnote{A semantic block describes a group of codes based on topical content. For example \textit{labelling}, \textit{sketching}, and \textit{annotating} were placed in the same semantic block.}. The full list of semantic blocks can be found in the Appendix. The resulting patterns were then reviewed manually for semantically similar codes that were not merged in the previous steps, producing a final set of 113 first order concepts.

    \item \textbf{Construction Gioia Data Structure:} The first order concepts were then grouped into second-order themes grouped by their subject, resulting in 17 second-order themes. The second-order themes were then organised into seven aggregate dimensions, reflecting each theme's place in the problem-solving process.
\end{enumerate}

The resulting seven aggregate dimensions are shortly described below, listing the contained second-order themes and highlighting representative verbs from the underlying pattern codes. The mapping is then visualised in Figure~\ref{fig:theme-dimension-mapping}. The complete pattern-code-to-theme mapping is provided in the appendix.

\begin{itemize}
    \item \textbf{Problem Cognition} is the largest dimension by statement count (537 statements), covering actions centred around performing maths and solving the problems. It comprises \textit{Problem Comprehension} (reading, restating, identifying), \textit{Solution Execution} (performing calculation, rearranging equation), and \textit{Non-LLM Offloading Channels} (taking notes, reviewing notes).
    \item \textbf{Metacognition} focuses on participants planning, evaluation, and correction of their own actions. The themes are \textit{Self-monitoring \& Evaluation} (verifying, evaluating, reflecting), \textit{Strategy Formulation} (planning, deciding, considering), and \textit{Errors \& Correction (Self)} (recognising, correcting).
    \item \textbf{LLM Reliance Decisions} covers decision-making about whether and how to rely on the \ac{LLM}. The themes \textit{Offloading Decisions} (deciding, considering) and \textit{Requests to LLM} (requesting information, requesting guidance) cover the two main aspects of this dimension.
    \item \textbf{Information Exchange} is centred around the communication between the participant and the \ac{LLM}. Its themes \textit{Query Formulation} (formulating), \textit{Information Disclosure to LLM} (submitting, stating), and \textit{LLM Response Processing} (reading, restating) cover the key steps of the interaction.
    \item \textbf{Mental Model \& Updates} groups themes focused on the participants' reception of \ac{LLM} outputs and the resulting attributions of mental states and capabilities. \textit{LLM Output Evaluation} (evaluating, scrutinising), \textit{LLM Guidance Uptake} (accepting, rejecting), \textit{Identifying \& Correcting LLM Errors} (correcting, criticizing), and \textit{LLM Attributions} (attributing).
    \item \textbf{Emotional Expression} is a single-theme dimension, covering pure emotional expression, such as frustration, surprise, satisfaction, and uncertainty.
    \item \textbf{Organisation} is constructed from a single theme, covering experiment logistics, tool handling, and clarification requests directed at the researcher rather than the system.
\end{itemize}

\begin{figure}[ht]
    \centering
    \resizebox{\textwidth}{!}{%
    \begin{tikzpicture}[
        verbbox/.style   = {neutralnode, align=left, text width=4.6cm, anchor=west,
                             font=\scriptsize, minimum height=0.5cm, inner sep=3pt},
        themelbl/.style  = {font=\footnotesize, align=left, text width=6.6cm, anchor=west},
        dimlbl/.style    = {font=\bfseries\small, text=paletteblue, align=left, text width=4.6cm, anchor=west},
        dimbox/.style    = {accentnode, inner sep=3pt},
        colheader/.style = {font=\bfseries\small, align=left, anchor=west},
    ]

        \pgfdeclarelayer{background}
        \pgfsetlayers{background,main}

        % fixed columns -- every node in a column uses this literal x, always. No exceptions.
        \def\colA{0}      % verb_box left edge
        \def\colB{5}      % theme label left edge
        \def\colC{12.3}   % dimension label left edge (frame width follows automatically via fit)

        % fixed rows -- y is the VERTICAL CENTER of each row (anchor=west), so group
        % pitch within a group: 0.6 | extra gap at a group boundary: 0.4

        % ============ Problem Cognition ============
        \node[verbbox]  (pc1-vb)  at (\colA,  0.0) {restating, extracting, identifying};
        \node[themelbl] (pc1-lbl) at (\colB,  0.0) {Problem Comprehension};
        \node[fit=(pc1-vb)(pc1-lbl), inner sep=0pt] (pc1) {};

        \node[verbbox]  (pc2-vb)  at (\colA, -0.6) {taking notes, reviewing notes};
        \node[themelbl] (pc2-lbl) at (\colB, -0.6) {Non-LLM Offloading Channels};
        \node[fit=(pc2-vb)(pc2-lbl), inner sep=0pt] (pc2) {};

        \node[verbbox]  (pc3-vb)  at (\colA, -1.2) {performing, restating, setting up};
        \node[themelbl] (pc3-lbl) at (\colB, -1.2) {Solution Execution};
        \node[fit=(pc3-vb)(pc3-lbl), inner sep=0pt] (pc3) {};

        \node[dimlbl] (dimPC-lbl) at (\colC, -0.6) {Problem Cognition}; % avg(0.00,-1.2)
        \begin{pgfonlayer}{background}
            \node[dimbox, fit=(pc1)(pc2)(pc3)(dimPC-lbl)] (dimPC) {};
        \end{pgfonlayer}

        % ============ Metacognition ============ (+0.4 group gap after pc3: -1.2-0.6-0.4)
        \node[verbbox]  (mc1-vb)  at (\colA, -2.2) {verifying, evaluating, reflecting};
        \node[themelbl] (mc1-lbl) at (\colB, -2.2) {Self-monitoring \& Evaluation};
        \node[fit=(mc1-vb)(mc1-lbl), inner sep=0pt] (mc1) {};

        \node[verbbox]  (mc2-vb)  at (\colA, -2.8) {planning, deciding, considering};
        \node[themelbl] (mc2-lbl) at (\colB, -2.8) {Strategy Formulation};
        \node[fit=(mc2-vb)(mc2-lbl), inner sep=0pt] (mc2) {};

        \node[verbbox]  (mc3-vb)  at (\colA, -3.4) {recognizing, correcting};
        \node[themelbl] (mc3-lbl) at (\colB, -3.4) {Errors \& Correction (Self)};
        \node[fit=(mc3-vb)(mc3-lbl), inner sep=0pt] (mc3) {};

        \node[dimlbl] (dimMC-lbl) at (\colC, -2.8) {Metacognition}; % avg(-2.2,-3.4)
        \begin{pgfonlayer}{background}
            \node[dimbox, fit=(mc1)(mc2)(mc3)(dimMC-lbl)] (dimMC) {};
        \end{pgfonlayer}

        % ============ LLM Reliance Decisions ============ (-3.4-0.6-0.4)
        \node[verbbox]  (rd1-vb)  at (\colA, -4.4) {requesting guidance, verification};
        \node[themelbl] (rd1-lbl) at (\colB, -4.4) {Requests to LLM};
        \node[fit=(rd1-vb)(rd1-lbl), inner sep=0pt] (rd1) {};

        \node[verbbox]  (rd2-vb)  at (\colA, -5.0) {deciding, considering};
        \node[themelbl] (rd2-lbl) at (\colB, -5.0) {Offloading Decisions};
        \node[fit=(rd2-vb)(rd2-lbl), inner sep=0pt] (rd2) {};

        \node[dimlbl] (dimRD-lbl) at (\colC, -4.7) {LLM Reliance Decisions}; % avg(-4.4,-5.0)
        \begin{pgfonlayer}{background}
            \node[dimbox, fit=(rd1)(rd2)(dimRD-lbl)] (dimRD) {};
        \end{pgfonlayer}

        % ============ Information Exchange ============ (-5.0-0.6-0.4)
        \node[verbbox]  (ie1-vb)  at (\colA, -6.0) {reading, restating};
        \node[themelbl] (ie1-lbl) at (\colB, -6.0) {LLM Response Processing};
        \node[fit=(ie1-vb)(ie1-lbl), inner sep=0pt] (ie1) {};

        \node[verbbox]  (ie2-vb)  at (\colA, -6.6) {formulating, planning};
        \node[themelbl] (ie2-lbl) at (\colB, -6.6) {Query Formulation};
        \node[fit=(ie2-vb)(ie2-lbl), inner sep=0pt] (ie2) {};

        \node[verbbox]  (ie3-vb)  at (\colA, -7.2) {submitting, stating};
        \node[themelbl] (ie3-lbl) at (\colB, -7.2) {Information Disclosure to LLM};
        \node[fit=(ie3-vb)(ie3-lbl), inner sep=0pt] (ie3) {};

        \node[dimlbl] (dimIE-lbl) at (\colC, -6.6) {Information Exchange}; % avg(-6.0,-6.6)
        \begin{pgfonlayer}{background}
            \node[dimbox, fit=(ie1)(ie2)(ie3)(dimIE-lbl)] (dimIE) {};
        \end{pgfonlayer}

        % ============ Mental Model & Updates ============ (-7.2-0.6-0.4)
        \node[verbbox]  (mm1-vb)  at (\colA, -8.2) {building, confirming, accepting};
        \node[themelbl] (mm1-lbl) at (\colB, -8.2) {LLM Guidance Uptake};
        \node[fit=(mm1-vb)(mm1-lbl), inner sep=0pt] (mm1) {};

        \node[verbbox]  (mm2-vb)  at (\colA, -8.8) {scrutinizing, evaluating};
        \node[themelbl] (mm2-lbl) at (\colB, -8.8) {LLM Output Evaluation};
        \node[fit=(mm2-vb)(mm2-lbl), inner sep=0pt] (mm2) {};

        \node[verbbox]  (mm3-vb)  at (\colA, -9.4) {correcting, criticizing};
        \node[themelbl] (mm3-lbl) at (\colB, -9.4) {Identifying \& Correcting LLM Errors};
        \node[fit=(mm3-vb)(mm3-lbl), inner sep=0pt] (mm3) {};

        \node[verbbox]  (mm4-vb)  at (\colA, -10.0) {anticipating, attributing};
        \node[themelbl] (mm4-lbl) at (\colB, -10.0) {LLM Attributions};
        \node[fit=(mm4-vb)(mm4-lbl), inner sep=0pt] (mm4) {};

        \node[dimlbl] (dimMM-lbl) at (\colC, -9.1) {Mental Model \& Updates}; % avg(-8.2,-10.0)
        \begin{pgfonlayer}{background}
            \node[dimbox, fit=(mm1)(mm2)(mm3)(mm4)(dimMM-lbl)] (dimMM) {};
        \end{pgfonlayer}

        % ============ Emotional Expression ============ (-10.0-0.6-0.4)
        \node[verbbox]  (ee1-vb)  at (\colA, -11.0) {expressing};
        \node[themelbl] (ee1-lbl) at (\colB, -11.0) {Emotional Expression};
        \node[fit=(ee1-vb)(ee1-lbl), inner sep=0pt] (ee1) {};

        \node[dimlbl] (dimEE-lbl) at (\colC, -11.0) {Emotional Expression};
        \begin{pgfonlayer}{background}
          \node[dimbox, fit=(ee1)(dimEE-lbl)] (dimEE) {};
        \end{pgfonlayer}

        % ============ Organisation ============ (-11.0-0.6-0.4)
        \node[verbbox]  (org1-vb)  at (\colA, -12.0) {managing logistics};
        \node[themelbl] (org1-lbl) at (\colB, -12.0) {Organisation};
        \node[fit=(org1-vb)(org1-lbl), inner sep=0pt] (org1) {};

        \node[dimlbl] (dimORG-lbl) at (\colC, -12.0) {Organisation};
        \begin{pgfonlayer}{background}
            \node[dimbox, fit=(org1)(dimORG-lbl)] (dimORG) {};
        \end{pgfonlayer}

        % ============ column headers ============
        \node[colheader] at (\colA, 0.8) {Representative Verbs};
        \node[colheader] at (\colB, 0.8) {Second-order Theme};
        \node[colheader] at (\colC, 0.8) {Aggregate Dimension};

        \end{tikzpicture}%
    }
    \caption[Second-order themes grouped into aggregate dimensions]{Second-order themes grouped into the seven aggregate dimensions, with representative verbs shown per theme. The complete grouping of pattern codes to aggregate dimensions can be found in the Appendix.}
    \label{fig:theme-dimension-mapping}
\end{figure}

It is important to note that dimensions have greatly varying sizes (with regard to the number of statements contained), and the size of a dimension is not a reliable indicator of its analytical importance. \textit{Problem Cognition} for example is the largest dimension, but its size is a result of the experimental setup that foregrounded the mathematics problem and allowed users to freely engage with the \ac{LLM}. The transcripts contained 1466 coded statements, with 888 statements from Group E, and 578 statements from group B, yielding global split of 60.6\% of statements from Group E. This split was used as a baseline to compare between-group differences in the aggregate dimensions, to determine whether the \ac{CoT} manipulation had an effect on the participants' behaviour. The statement counts per aggregate dimension are shown in Table~\ref{tab:statement-counts}.

\begin{ctable}
    \centering
    \begin{tabularx}{\textwidth}{X l l l}
        \toprule
        \textbf{Aggregate Dimension} & \textbf{Statements} & \textbf{Count (E)} & \textbf{Count (B)} \\
        \midrule
        Problem Cognition & 537 & 341 (63.5\%) & 196 (36.5\%)\\
        Metacognition & 325 & 193 (59.4\%) & 132 (40.6\%)\\
        Information Exchange & 248 & 167 (67.3\%) & 81 (32.7\%)\\
        LLM Reliance Decisions & 121 & 80 (66.1\%) & 41 (33.9\%)\\
        Mental Model \& Updates & 115 & 56 (48.7\%) & 59 (51.3\%)\\
        Organisation & 66 & 28 (42.4\%) & 38 (57.6\%)\\
        Emotional Expression & 54 & 23 (42.6\%) & 31 (57.4\%) \\
        \midrule
        \textbf{Total} & \textbf{1466} & \textbf{888 (60.6\%)} & \textbf{578 (39.4\%)} \\
        \bottomrule
    \end{tabularx}
    \caption[Statement counts per aggregate dimension]{Statement counts per aggregate dimension, with the number of statements from each group and the corresponding share in parentheses.}
    \label{tab:statement-counts}
\end{ctable}

\textit{Organisation} serves as a control for the between-group comparison. Since it is composed of questions surrounding the experiment and material supplied to the participants, it is not plausibly sensitive to the \ac{CoT} manipulation. Consequently, any deviation from the baseline smaller than the deviation of \textit{Organisation} should not be used as sole evidence for a between-group difference. Since \textit{Organisation} shows the largest deviation (-18.2\%) from the baseline, dimension-level between-group comparisons cannot be used to sensibly derive design principles. Conversely, \textit{Organisation} is among the smallest dimensions. Notably, the only dimension with a similar deviation from the baseline is \textit{Emotional Expression} (-18.0\%), indicating that the small dimension size could play a role in this observation. For this reason, the between-group comparisons are not used to derive design principles, but rather as supporting evidence for the qualitative findings. As a result, the empirically derived \acp{DP} do not address the explanations as a mechanism. Section~\ref{ssec:uncovered-design-space} discusses this gap in the design space and proposes theory-based design principles to address it.

Comparisons of behavioural patterns across the two groups are however still possible. This applies to the finding that Group B participants are more likely to heuristically accept the \ac{LLM}'s guidance than Group E participants. This was established by examining the \textit{Mental Model \& Updates} dimension, which contains three second-order themes that are directly relevant to the appraisal of the \ac{LLM}'s output: \textit{LLM Guidance Uptake}, \textit{LLM Output Evaluation}, and \textit{Identifying \& Correcting LLM Errors}. \textit{LLM Attributions} is excluded because it is not directly relevant to the appraisal of the \ac{LLM}'s output, but rather to the participant's understanding of the chatbot's capabilities. The three appraisal themes contain 96 rows (46 E, 50 B). Within this subset, four pattern codes were classed as heuristic uptake: \textit{confirming LLM's guidance}, \textit{accepting LLM's assertions}, \textit{deciding to follow LLM guidance}, and \textit{following LLM guidance}; the remainder was classified as deliberation. Calculating the share of heuristic uptake within each group yields 17.4\% for Group E (8/46) against 36.0\% for Group B (18/50). This finding is not undermined by the \textit{Organisation} control, because the ratios are constructed purely within-group, so that any difference in verbosity between groups cancels out. To further validate the finding, the pool of \textit{reading LLM's guidance} and \textit{restating LLM's guidance} statements was re-examined for heuristic acceptances absorbed into those two codes. As a result of this re-examination, one row was reassigned into the heuristic-uptake count reported above, and two additional rows were assigned to other codes. To check the stability of the finding, results were validated using \ac{LOO} validation, which showed that Group B's heuristic share is sensitive to individual participants (specifically P11), but the direction of the finding is stable.

The same subset also permits an observation about the consequences of appraisal. Of the 96 \ac{LLM}-directed appraisal statements, only 7 (about 7\%) resulted in a correction of the system. A mirrored set can be constructed on the participant side, pairing a scrutiny theme with a correction theme in the same way: \textit{Self-monitoring \& Evaluation} (137) and \textit{Errors \& Correction (Self)} (69) yield 206 self-directed appraisal statements, of which roughly 33\% resulted in a self-correction. Because both sets are constructed on the same principle and each ratio is calculated within its own set, the comparison is unaffected by the differing volume of statements on each side. The imbalance of corrections based on appraisal is taken up in Section~\ref{ssec:discussion-results}.

From the data structure presented above, a total of six \acp{DP} were derived using, by looking for codes, themes and dimensions that directly serve or violate the \acp{MR} and \acp{DR}. Additionally, the data structure was examined for patterns showing direct attribution of mental states to the \ac{LLM}. \acp{DP} address how a \ac{DR} can be realised in a design theory \parencite{Herm2022}. The \acp{DP} are presented below, including a rationale, and evidence used to derive and validate the principle. Select quotes from the transcripts are used where appropriate and when they are representative of the underlying behavioural pattern. Each quote is referenced by its unique identifier, as described above after the translation of the original German quote into English. The initial \acp{DP} are denoted by the letters A-F, to avoid confusion with final \acp{DP} numbering presented in Section~\ref{sssec:revised-dp}.

\begin{enumerate}[label=\textbf{DP-\Alph*}, leftmargin=*]
    % former DP-F
    \item \textbf{Re-use the user's own task vocabulary}, proposes to adopt task-relevant terminology that the user has introduced. It is based on the observation, that users mapped repeated information from the \ac{LLM} onto their own vocabulary. The relevant pattern code \textit{reading LLM's restated known information} contains 29 instances well distributed across both groups. This behaviour also aligns with \textcite{Shin2021}'s account of user-relative intelligibility (cf. Section~\ref{ssec:explanations-mental-model-formation}). The \ac{DP} therefore serves \textbf{DR1}, by making the system's actions more comprehensible to each individual user.

    % former DP-G
    \item \textbf{Clearly state system limitations}, to the user during the conversation. This \ac{DP} is motivated by the observation that participants attributed capabilities to the system, without factual basis (pattern code \textit{attributing capability to LLM}). A good example is the expectation, that the \ac{LLM} draws the same conclusion from the provided image (demonstrating an implicit belief attribution): \enquote{aber eigentlich hat er ja das Bild} [but he does have the image] (ID: \texttt{11.014.1}). The attempts to get the \ac{LLM} to solve the maths problems are a prominent example, where the system should clearly state its limitations. By clearly communicating limitations, the user expectations about system capabilities are managed fulfilling \textbf{DR2} \parencite[cf.][]{Amershi2019}.

    % former DP-H
    \item \textbf{Support error detection}, focuses on supporting the user in detecting and ascribe errors. This \ac{DP} is anchored in the observation, that users accepted system advice heuristically (as evidenced by the share of heuristic uptake described above). Highlighting the source (e.g., task given, user-provided, externally-retrieved, etc.) of a piece of information in the system response could invite more scrutiny by the users. Additionally, the \ac{LLM} also accepted user input, even if it conflicted with previous inputs. The most prominent example of this behaviour is, when the system adopted a sign error by P08 signalling acceptance of the change \parencite{Clark1989}. The participant later explicitly commented \enquote{Schade, dass er solche Logikfehler nicht erkennt.} [Shame he doesn't recognise such logical errors.] (ID: \texttt{08.048.6}). This design principle serves to fulfil \textbf{DR1} and \textbf{DR3} by providing sources for information \parencite[cf.][]{Steyvers2025}, and \textbf{DR4} by surfacing conflicting information to the user.

    % former DP-C
    \item \textbf{Make system assumptions inspectable and correctable}, is founded on the observation that participants frequently detected (perceived) errors in the system's assumptions. The data for \textit{Mental Model \& Updates} shows that participants validate the system's outputs by seeking and validating assumptions in the response, demonstrating belief attributions focusing on misrepresentation \parencite{Rakoczy2022}. The most prominent example is a participant stating \enquote{Deine Annahme \ldots ist falsch} [Your assumption \ldots is wrong] (ID: \texttt{08.016.1}) in a message to the \ac{LLM}, followed by a correction. This \ac{DP} also leans on the finding that Group E participants are more likely to deliberate on the system's outputs than Group B participants, as evidenced by the lower heuristic uptake share (17.4\% against 36.0\%). Similar to \ac{CoT} explanations, surfacing assumptions exposes the system's internal state, which is then deliberated on by the user. The mechanism seeks to prevent the failure mode \textit{\ac{AI}-based confirmation} by supporting conscious deliberation \parencite{Jussupow2021}. This \ac{DP} serves \textbf{DR1}, by making the system's actions more comprehensible \parencite[cf.][]{Amershi2019}, and \textbf{DR5}, by reducing the effort otherwise required to infer why the system behaved as it did.

    % former DP-A
    \item \textbf{Establish request purpose before querying}, aims to force users to deliberately consider what kind of support they are expecting from the system before issuing a request. This \ac{DP} is primarily motivated by two related failure patterns, in which a lack of established expectations or purpose leads to dissatisfaction with the system's response. In the first pattern, participants did not articulate a role at all, to then express their dissatisfaction with the system's response in retrospect, e.g., \enquote{hat es mir jetzt nicht so viel weiter geholfen} [It didn't help me much] (ID: \texttt{10.007.1}). In a related pattern, participants explicitly articulated a query that was incompatible with the system's instructions, which instructed it explicitly not to solve the maths problems (cf. \textbf{R3}, Section~\ref{ssec:chatbot-implementation}), e.g., \enquote{löse es [Gleichungssystem] mir auf und berechne die lösungen} [solve it for me and calculate the solutions] (ID: \texttt{03.024.1}). Both failure patterns can be explained by a lack of consideration (and communication), what is expected from the system. This is also supported by quantitative data: Of the 40 total codes in which participants committed to querying the system (\textit{deciding to ask LLM for guidance/information/verification/solution}), in 17 (42.5\%) the role was inferred from the statements' context rather than explicitly articulated. The principle is intended to address this failure by forcing users to consider and articulate the purpose/goal of their request before issuing it. This \ac{DP} fulfils \textbf{DR4}, by introducing friction at a consequential decision \parencite[cf.][]{Buinca2021, Weisz2025}. It also fulfils \textbf{DR2}, as establishing a role sets a corresponding expectation of behaviour and capability \parencite[cf.][]{Bansal2019}. When applied correctly, the \ac{DP} also helps to reduce cognitive load (\textbf{DR5}), as the improved alignment between request and response reduces the need for further deliberation.

    % former DP-E
    \item \textbf{Provide additional, user-owned working surfaces}, seeks to add additional interaction surfaces alongside the conversational interface, on which users can externalise and manipulate intermediate work. This \ac{DP} is motivated by the observation that participants frequently took notes. Depending on the specific use case, the conversational interface could even be provided as a secondary feature in a larger application, similar to how coding assistants like GitHub Copilot \parencite{GitHubInc2026} or Claude Code \parencite{AnthropicInc2026} are integrated into IDEs. This also aligns with a well documented behaviour called cognitive offloading \parencite{Risko2016}, which describes the tendency to manipulate the environment to support cognitive processes. Participants explicitly stated, they took notes to support their cognitive processes: \enquote{weil ich kann besser denken, wenn ich die Sachen tatsächlich daran schreibe} [because I can think better when I actually write things down] (ID: \texttt{05.025.2}). An \ac{LLM}-readable surface allows the system to access and incorporate the user's intermediate work, reducing the need for manual tracking and re-statement of the information, helping to fulfil \textbf{DR5}.
\end{enumerate}

Across the six design principles, coverage of the five design requirements varies both in how many principles address a given \ac{DR} and in the mechanism through which they do so. \textbf{DR1} receives the broadest support, addressed by \textbf{DP-A}, \textbf{DP-C}, and \textbf{DP-D}, each through a distinct category of content: task vocabulary, information provenance, and system assumptions, respectively. \textbf{DR2} is addressed by \textbf{DP-B} and \textbf{DP-E}, which act at opposite points in the interaction: \textbf{DP-B} aligns expectations through the system's own disclosure during the response, while \textbf{DP-E} aligns them implicitly before a request is even issued. \textbf{DR4} is addressed by \textbf{DP-C} and \textbf{DP-E}, and \textbf{DR5} by \textbf{DP-D}, \textbf{DP-F}, and, to a lesser extent, \textbf{DP-E}. \textbf{DR3} stands apart, as the only design requirement addressed by a single principle, \textbf{DP-C}; this thin coverage is picked up again in Section~\ref{ssec:uncovered-design-space}, alongside the explanation-mechanism gap. Figure~\ref{fig:mr-dr-dp-mapping} visualises the complete mapping between meta-requirements, design requirements, and the six initial design principles.

\begin{figure}[ht]
    \centering
    \resizebox{\textwidth}{!}{%
    \begin{tikzpicture}[
        node distance=0pt,
        mrbox/.style={neutralnode, align=left, text width=5.1cm},
        drbox/.style={neutralnode, align=left, text width=6cm},
        dpbox/.style={neutralnode, align=left, text width=5.5cm},
        mr1line/.style={paletteblue, line width=1pt},
        mr2line/.style={palettered, line width=1pt},
        mr3line/.style={paletteochre, line width=1pt},
        dr1line/.style={palettegreen, line width=1pt},
        dr2line/.style={paletteochre, line width=1pt},
        dr3line/.style={paletteviolet, line width=1pt},
        dr4line/.style={paletteblue, line width=1pt},
        dr5line/.style={palettered, line width=1pt},
        partial/.style={densely dotted},
    ]
        % ---- Meta-requirements ----
        \node[mrbox] (mr1) at (0,2.3)  {\textbf{MR1:}~Support accurate mental models};
        \node[mrbox] (mr2) at (0,0)    {\textbf{MR2:}~Enable appropriate reliance};
        \node[mrbox] (mr3) at (0,-2.3) {\textbf{MR3:}~Support appropriate cognitive engagement};

        % ---- Design requirements ----
        \node[drbox] (dr1) at (8.5,3.0)  {\textbf{DR1:}~Increase intelligibility of system behaviour};
        \node[drbox] (dr2) at (8.5,1.5)  {\textbf{DR2:}~Align capability expectations with actual system capability};
        \node[drbox] (dr3) at (8.5,0.0)  {\textbf{DR3:}~Increase verifiability of system outputs};
        \node[drbox] (dr4) at (8.5,-1.5) {\textbf{DR4:}~Increase deliberate engagement at consequential decision points};
        \node[drbox] (dr5) at (8.5,-3.0) {\textbf{DR5:}~Decrease cognitive effort of interpreting system outputs};

        % ---- Design principles ----
        \node[dpbox] (dp1) at (17,3.5)  {\textbf{DP-A:}~Re-use the user's own task vocabulary};
        \node[dpbox] (dp2) at (17,2.0)  {\textbf{DP-B:}~Clearly state system limitations};
        \node[dpbox] (dp3) at (17,0.75)  {\textbf{DP-C:}~Support error detection};
        \node[dpbox] (dp4) at (17,-0.5) {\textbf{DP-D:}~Make system assumptions inspectable and correctable};
        \node[dpbox] (dp5) at (17,-2) {\textbf{DP-E:}~Establish request purpose before querying};
        \node[dpbox] (dp6) at (17,-3.5) {\textbf{DP-F:}~Provide additional, user-owned working surfaces};

        % ---- Column headings ----
        \node[font=\bfseries] at (0,4.7) {Meta-Requirements};
        \node[font=\bfseries] at (8.5,4.7) {Design Requirements};
        \node[font=\bfseries] at (17,4.7) {Design Principles};

        % ---- MR -> DR edges, coloured by source MR (cf. Figure~\ref{fig:mr-dr-mapping}) ----
        \draw[mr1line]          (mr1.east) -- (dr1.west);
        \draw[mr1line]          (mr1.east) -- (dr2.west);
        \draw[mr2line]          (mr2.east) -- (dr1.west);
        \draw[mr2line]          (mr2.east) -- (dr2.west);
        \draw[mr2line]          (mr2.east) -- (dr3.west);
        \draw[mr2line]          (mr2.east) -- (dr4.west);
        \draw[mr3line]          (mr3.east) -- (dr1.west);
        \draw[mr3line,partial]  (mr3.east) -- (dr3.west);
        \draw[mr3line]          (mr3.east) -- (dr4.west);
        \draw[mr3line]          (mr3.east) -- (dr5.west);

        % ---- DR -> DP edges, coloured by source DR; dotted where DP prose marks the service as partial ----
        \draw[dr1line]          (dr1.east) -- (dp1.west);
        \draw[dr2line]          (dr2.east) -- (dp2.west);
        \draw[dr1line]          (dr1.east) -- (dp3.west);
        \draw[dr3line]          (dr3.east) -- (dp3.west);
        \draw[dr4line]          (dr4.east) -- (dp3.west);
        \draw[dr1line]          (dr1.east) -- (dp4.west);
        \draw[dr5line]          (dr5.east) -- (dp4.west);
        \draw[dr2line]          (dr2.east) -- (dp5.west);
        \draw[dr4line]          (dr4.east) -- (dp5.west);
        \draw[dr5line, partial] (dr5.east) -- (dp5.west);
        \draw[dr5line]          (dr5.east) -- (dp6.west);
    \end{tikzpicture}%
    }
    \caption[Meta-requirements, design requirements, and initial design principles]{The initial design principles (pre-revision) mapped onto the design requirements which they fulfil, extending Figure~\ref{fig:mr-dr-mapping}.}
    \label{fig:mr-dr-dp-mapping}
\end{figure}

\subsection{Design Principle Evaluation \& Revision} \label{ssec:design-principle-revision}

This section covers the evaluation of the \acp{DP} and the revision that resulted from the evaluation. It is split into Section~\ref{sssec:interview-procedure} describing the interview procedure and experts that participated in the interview. In Section~\ref{sssec:interview-results} the key results of the evaluation are presented, and Section~\ref{sssec:revised-dp} presents the revised, final set of \acp{DP}.

\subsubsection{Interview Procedure and Participants} \label{sssec:interview-procedure}

The six design principles derived in Section~\ref{ssec:initial-design-principles} were evaluated through semi-structured expert interviews. In total four experts were interviewed, three of whom were \ac{HCI} and \ac{HAII} researchers, and one industry practitioner with focus on the development of generative AI applications. Details about all participants are listed in Table~\ref{tab:expert-participants}.

\begin{ctable}
    \centering
    \begin{tabular}{c c c c}
        \toprule
        \textbf{ID} & \textbf{Education} & \textbf{Current Role} & \textbf{Experience} \\
        \midrule
        E01 & M.Sc. (\acs{HCI}) & PhD candidate (\acs{HAII}) & 8 months \\
        E02 & M.Sc. (CS) & PhD candidate (\acs{HCI}) & 8 months \\
        E03 & B.Sc. (CS) & Software Engineer \& Designer & 5 years \\
        E04 & M.Sc. (\acs{HCI}) & PhD candidate (\acs{HAII}) & 6 months \\
        \bottomrule
    \end{tabular}
    \caption[Expert interview participants]{Overview of the four experts, including highest degree of education, current professional role, and years of experience in their current role.}
    \label{tab:expert-participants}
\end{ctable}

The semi-structure format allowed for the identical evaluation of the \ac{DP} across all experts, while also allowing for the exploration of additional, unprompted topics. The interviews were conducted online via Zoom, and lasted between 60 and 75 minutes and followed the same high level procedure outlined below. The interviews were recorded and transcribed for analysis.

\begin{enumerate}
    \item \textbf{Consent and Demographics:} Participants were informed about the goal of the interview, the use of the collected data, and that the interview would be recorded. After consent was obtained, demographic information was collected to contextualise the expert's background and experience.
    \item \textbf{Topic Introduction:} The thesis topic and relevant concepts (\ac{ToM}, \ac{DSR}, and \acp{DP}) were introduced to the participants as needed for their individual background. The thesis methodology was explained indiscriminately to all participants.
    \item \textbf{Comprehension Check:} The \acp{MR} and \acp{DR} were presented, and their relationships explained. Before evaluating each \ac{DP} a comprehension check was performed using a fictional example of the same \ac{DR}-mapping exercise used later for the \acp{DP}. This validated the expert's understanding of the \ac{MR}/\ac{DR}/\ac{DP} concepts and familiarised them with the exercise format ahead of the main discussion. All four experts passed the check.
    \item \textbf{\ac{DP} Evaluation:} Each \ac{DP} was treated identically, with a predefined set of questions asked. The \ac{DP} was first presented to the expert, followed by a discussion of its \ac{DR}-mapping and three questions, targeting \textbf{actability and guidance}, \textbf{importance}, and \textbf{novelty and insightfulness} \parencite{Iivari2021}.
    \item \textbf{Closing Discussion:} After all \acp{DP} were discussed, a final set of questions was asked to evaluate the complete set of \acp{DP} to assess their \textit{effectiveness} \parencite{Iivari2021} and \textit{deficit} \parencite{Janiesch2020}. A question to assess \textit{redundancy} and \textit{excess} was planned, but omitted due to an oversight in the interview guide. The interview was closed with two open questions to allow the expert to provide any additional feedback.
\end{enumerate}

\textbf{Accessibility} \parencite[cf.][]{Iivari2021} was not assessed using a dedicated question, but instead inferred from the discussion of each \ac{DP} and the expert's \ac{DR}-mapping. Two key indicators were used to operationalise accessibility: Firstly, a qualitative measure, based on the clarification-seeking questions raised by the expert during the discussion (e.g., requesting the meaning of a term or an example). Secondly, a quantitative measure, based on the agreement between the expert's stated \ac{DR}-mapping and the mapping intended by the derivation in Section~\ref{ssec:initial-design-principles}. The criterion of \textbf{overload} \parencite[cf.][]{Janiesch2020} was not assessed directly, but an indirect signal can be inferred from the actability question, based on the variance in experts' understanding of the mechanism. The complete set of questions asked during the interviews is provided in the Appendix.

\subsubsection{Result Summary} \label{sssec:interview-results}

The interview transcripts were evaluated after the interviews by filling an evaluation form recording the \ac{DR}-mappings and responses to each question. The responses to each question were then combined into an evaluation based on the categories described in \ref{sssec:interview-procedure}. The \ac{DR}-mappings are shown in Table~\ref{tab:dr-mapping-experts} compared to the mapping proposed in Section~\ref{ssec:initial-design-principles}. Individual remarks beyond this fixed set of categories are included below only where they were raised independently by two or more experts, or where they connect directly to a construct or citation already established elsewhere in the design theory. The qualitative evaluation of each \ac{DP} is outlined below, followed by notions from the closing discussion.

\begin{ctable}
    \centering
    % LTeX: enabled=false
    \begin{tabular}{c c c c c c}
        \toprule
        \textbf{DP} & \textbf{Reference} & \textbf{E01} & \textbf{E02} & \textbf{E03} & \textbf{E04} \\
        \midrule
        \textbf{DP-A} & DR1 & DR1, \textcolor{paletteochre}{DR5} & DR1, \textcolor{paletteochre}{DR5} & DR1 & \textcolor{palettered}{DR2} \\
        \textbf{DP-B} & DR2 & DR2* & \textcolor{paletteochre}{DR1}, DR2 & DR2* & DR2 \\
        \textbf{DP-C} & DR1, DR3, DR4 & DR3, DR4 & DR3 & DR3 & DR3 \\
        \textbf{DP-D} & DR1, DR5 & \textcolor{palettered}{DR3} & \textcolor{palettered}{DR3} & \textcolor{palettered}{DR3} & \textcolor{palettered}{DR3} \\
        \textbf{DP-E} & DR2, DR4, (DR5) & DR2 & DR2 & \textcolor{palettered}{DR1} & \textcolor{paletteochre}{DR1}*, DR2 \\
        \textbf{DP-F} & DR5 & \textcolor{paletteochre}{DR4}* or DR5* & \textcolor{palettered}{DR4} & DR5* & DR5* \\
        \bottomrule
    \end{tabular}
    % LTeX: enabled=true
    \caption[Expert \ac{DR}-mapping of the initial design principles]{Expert-stated \ac{DR}-mapping for each initial \ac{DP}, against the mapping shown in Section~\ref{ssec:initial-design-principles}. \textcolor{paletteochre}{Yellow} marks an added \ac{DR}, \textcolor{palettered}{red} marks full deviation, and an asterisk marks a mapping voiced only after clarification.}
    \label{tab:dr-mapping-experts}
\end{ctable}

The results of the evaluation show that the mappings of the experts are generally in agreement with the proposed mappings. The most notable differences are a complete misalignment of \textbf{DP-D} and the high amount of clarification required for \textbf{DP-F}. In contrast, \textbf{DP-A}, \textbf{DP-B}, and \textbf{DP-C} are well understood, with \textbf{DP-E} showing minor comprehension challenges. Another notable difference between the experts mappings and the proposed mappings is that the experts primarily mapped only one \ac{DR} to each \ac{DP}, indicating that the experts' understanding of the \acp{DP} was more focused than the proposed mappings.

In the following the results for each \ac{DP} are discussed in detail, including the \ac{DR}-mapping. Relevant remarks from the interviews will be highlighted, and the implications for the design principles will be discussed.

\begin{enumerate}[label=\textbf{DP-\Alph*}, leftmargin=*]
    % former DP-F
    \item \textbf{Re-use the user's own task vocabulary} \\
        Three of the four experts mapped the \ac{DP} to \textbf{DR1}, with two adding \textbf{DR5} alongside it. E04 was the outlier, mapping the \ac{DP} to \textbf{DR2}, arguing that assigning a role creates user expectations regarding the system's capabilities. This shows the \ac{DP} is very \textit{accessible}. The \ac{DP}-specific discussion questions showed, experts found the \ac{DP} to be \textit{actable} and \textit{important}, but not particularly \textit{novel}, calling it existing behaviour. During the discussion, E01 and E04 raised the risks of sycophancy and anthropomorphism, respectively. They warned that mirroring the user's vocabulary could lead to reduced scrutiny and an increased tendency to anthropomorphise the system. These risks are also raised in the literature \parencite{Sharma2024, Colombatto2025}, leading to the conclusion that the \ac{DP} should be revised to mitigate these risks.

    % former DP-G
    \item \textbf{Clearly state system limitations} \\
        All four experts agreed, that \textbf{DP-B} addresses \textbf{DR2}, with the caveat that two required clarifications, before being able to map the \ac{DP}. This speaks to a good \textit{accessibility}, with improvement potential in the mechanism, which caused both clarification questions. Experts agreed, that the \textit{novelty} of the \ac{DP} was limited. Experts suggestions for implementation almost all shared the idea of tweaking the vocabulary the model uses in its responses, indicating clear \textit{actability}. It will be revised by rewording the mechanism of the \ac{DP} to focus on the system's own disclosure of its limitations.
        % This sentence belongs in the discussion; contrast deliberate choice for analogy with anthropophormism risk
        % During the open discussion, E03 raised that while the use of certain verbs (e.g., \enquote{thinking} or \enquote{reasoning}) invite anthropomorphism, they can be also understood as analogies.

    % former DP-H
    \item \textbf{Support error detection} \\
        The \textit{actability} question revealed two directions for \textbf{DP-C}, revealing two distinct mechanisms to support error detection. The first one focused on providing sources for information, through external sources or references to the existing chat history. The second mechanism focused on surfacing conflicts between established context and new information provided by the user. Combined with the \ac{DR}-mapping showing high \textit{accessibility}, this indicated that the \ac{DP} was \textit{overloaded} and should be split into two separate principles. All experts agreed that the \ac{DP} was \textit{important}, with E01 noting that instructing the system to surface conflicts could help counteract the sycophancy risk raised in \textbf{DP-A}. Similarly, E03 noted that surfacing conflicting info might contradict the sycophantic tendencies of \acp{LLM}. \textit{Novelty} was rated high by all experts as well, with all experts noting that combining external and internal references was a new idea to them.

    % former DP-C
    \item \textbf{Make system assumptions inspectable and correctable} \\
        All experts mapped \textbf{DP-D} to \textbf{DR3}, showing uniform disagreement with the proposed mapping. The uniformity of the mapping and the lack of clarifying questions, suggest that this is not an \textit{accessibility} problem. Instead, the \ac{DR}-mapping of this \ac{DP} will be reevaluated. Implementation concepts converged on the idea of highlighting assumptions in the system's response using uncertainty as a determinant, e.g., by using hedging language or confidence scores. This indicates that the \ac{DP} is \textit{actable}. It is however noteworthy, that the idea of self-reporting confidence is currently contested in literature \parencite{Kapoor2024}.

    % former DP-A
    \item \textbf{Establish request purpose before querying} \\
        \textbf{DP-E} was mapped to \textbf{DR2} by most experts, indicating average \textit{accessibility}. A notable omission by all experts is a lack of mapping to \textbf{DR4}, as E01 and E04 both raised concerns about the friction introduced by the \ac{DP}, which could lead to habituation and reduced effectiveness over time. The \textit{importance} was rated high by all experts, with E01 noting very specialised systems (e.g., interactive knowledge databases) as exceptions, where the \ac{DP} would be less important. Three of the experts said that while the underlying idea of tailoring the system's response to a given role was known to them, the specific mechanism of enforcing a role assignment was genuinely new.

    % former DP-E
    \item \textbf{Provide additional, user-owned working surfaces} \\
        \textbf{DP-F} showed the lowest \textit{accessibility} of all \acp{DP}, with three of four experts requiring clarification before they could map the \ac{DP}. The biggest challenge was the vague wording of the \ac{DP}, leading to uncertainty about the intended mechanism. After clarification, experts rated the \textit{novelty} of the \ac{DP} as low, despite E01 noting, that the notion that system readable working surfaces reduce cognitive load was new to them. This indicates a need to reword the \ac{DP} using simpler language.
\end{enumerate}

% TODO: find canonical source for automation bias
During the mapping exercise, E01, E03, and E04 all asked about the term \enquote{intelligibility} and noted that it could be replaced with \enquote{understandability}. The complete set of \acp{DP} was evaluated in the categories of \textit{effectiveness} \parencite{Iivari2021} and \textit{deficit} \parencite{Janiesch2020}. As the question for \textit{redundancy} and \textit{excess} was omitted and experts did not voice any unprompted comments, no evaluation could be performed in these categories. The experts agreed, that the set of \acp{DP} would help to \textit{effectively} address the \acp{DR}, helping to improve the quality of systems built with the \acp{DP}. Experts also agreed, that the language of the \acp{DP} and doubts about full technical feasibility of the mechanisms were the main challenges to overcome when using the \acp{DP}. When asked about \textit{deficits}, E01 noted that the \ac{DP} could be complemented addressing known cognitive biases (e.g., automation bias). While automation bias \parencite{PLACEHOLDER} in itself is not covered in this thesis, the associated risk of heuristic acceptance \parencite{Jussupow2021} is discussed in Section~\ref{ssec:mental-state-attributions}. E02 was the only expert to ask why none of the \acp{DP} addressed the \ac{CoT} explanations directly, which will be addressed in Section~\ref{ssec:uncovered-design-space}.

\subsubsection{Revised Design Principles} \label{sssec:revised-dp}

The findings from Section~\ref{sssec:interview-results} were used to revise the existing \ac{DP}. The final DP are presented in this section using the reduced Gregor format described in Section~\ref{ssec:thesis-methodology}. Specifying the \textbf{context} of the DP, is omitted because the context is identically scoped to a collaborative problem-solving setting, where a user is assisted by an \ac{LLM}. Each final DP is presented below in the same format, with the title at the top followed \textbf{implementer}, \textbf{aim}, \textbf{user}, \textbf{mechanism}, and \textbf{rationale} in tabular form.

\textbf{DP-A} required substantial revision, to address experts concerns about anthropomorphic tendencies, when mirroring user vocabulary. The revised \textbf{DP1} broadens the scope from vocabulary-matching to adjusting the system's response behaviour to the user's demonstrated (or stated) expertise level. Broadening the mechanism allows the mechanism to still address user-specific comprehension \parencite{Shin2021}, while avoiding prescribing potentially sycophantic behaviour \parencite{Sharma2024} and anthropomorphic tendencies \parencite[cf.][]{Weisz2025} raised by the experts. The experiment transcripts provide an example of the challenges introduced, when the system's responses are not aligned with the user's expertise level (cf. ID \texttt{11.009.1} and \texttt{11.017.1}ff). The rationale for the revised \textbf{DP1} can be retained. The revised \textbf{DP1} continues to serve \textbf{DR1} and now also addresses \textbf{DR5}, as matching the user's expertise level reduces the effort required especially for users with lower expertise, who would otherwise need to infer the meaning of technical terms and concepts.

\designprinciple{DP1}{Adapt responses to user expertise}
    {For the \textit{developer} to adapt the \textit{system's responses} to the \textit{user's demonstrated expertise level}}
    {use the \textit{user's vocabulary and wording} as an initial signal of expertise, and \textit{adapt the complexity, technicality, and level of depth in subsequent responses} to match the user's demonstrated expertise}
    {Because matching the user's expertise increases the understandability of the system's responses.}

\textbf{DP-B} is functionally unchanged, but the mechanism is reworded, to improve the \textit{accessibility} of the \ac{DP}. The revised mechanism shifts the emphasis from the user's inference of system capabilities to the system's own disclosure of its limitations. The rationale is unchanged, as managing user expectations about system capabilities reduces the likelihood of over- or under-reliance \parencite[cf.][]{Jussupow2021}.

\designprinciple{DP2}{Disclose system limitations}
    {For the \textit{developer} to tune the system to \textit{disclose its limitations} to \textit{users} during the interaction}
    {\textit{communicate the system's limitations} clearly when they become relevant to the current task}
    {Because not disclosing system limitations can lead to incorrect capabilities being inferred by the user.}

\textbf{DP-C} is split into two separate principles to address the two distinct mechanisms identified through the interviews. The new \textbf{DP3} focuses on the information provenance mechanism, that provides sources for information, to help the user estimate information accuracy and support validation \parencite[cf.][]{Steyvers2025}. \textbf{DP4} prescribes the surfacing of conflicts between established context and new user input, to prevent silent endorsement by the \ac{LLM} \parencite{Clark1989}. \textbf{DP4} still carries the benefits of reducing the sycophancy risk noted by the experts. The \ac{DR}-mapping of the new principles still covers the same \acp{DR} as the original \textbf{DP-C}, but now with a more focused mechanism for each principle. \textbf{DP3} now helps fulfil \textbf{DR1} and \textbf{DR3}, as providing sources for information, allows users to better understand the system's use of information. \textbf{DP4} serves \textbf{DR4}, as surfacing conflicts between established context and new user input increases deliberate engagement.

\designprinciple{DP3}{Indicate information provenance}
    {For the \textit{developer} to \textit{increase the understandability of the system} for \textit{users}}
    {provide the \textit{sources of information} in the system's response (e.g., user-provided, system-assumed, or externally retrieved).}
    {Because indicating the source of information invites scrutiny of system-attributed claims rather than heuristic acceptance, increasing both the understandability of the system's use of information and the verifiability of individual claims}

\designprinciple{DP4}{Surface conflicting information}
    {For the \textit{developer} to \textit{support the detection of user errors} for \textit{users}}
    {\textit{surface conflicts} between established information from the context and new user input.}
    {Because silently adopting contradicting input signals endorsement instead of inviting deliberation. Surfacing the conflict instead invites the user to deliberate whether the change is intended and warranted.}

\textbf{DP-D} is kept as initially proposed, but the \ac{DR}-mapping is revised to reflect the experts' unanimous mapping to \textbf{DR3}. Because the mechanism and rationale remain unchanged, \textbf{DR3} is added to the mapping alongside the original \textbf{DR1} and \textbf{DR5}. The mechanism and rationale are unchanged, as they elicited very high \textit{importance} ratings from all experts.

\designprinciple{DP5}{Make system assumptions inspectable and correctable}
    {For the \textit{developer} to make the \textit{assumptions inspectable and correctable} for \textit{users}}
    {provide an editable overview of the \textit{assumptions underlying the system's response}.}
    {Because surfacing assumptions enables users to validate the system's intermediate steps, increasing the verifiability and therefore also increasing the understandability of system behaviour. Making the assumptions editable allows users to easily correct them, reducing cognitive effort.}

Conceptually, \textbf{DP-E} remains unchanged, but the mechanism is reworded to improve the \textit{accessibility} of the \ac{DP}. The revised mechanism focuses on establishing the collaboration mode and the revised rationale addresses the intentionality of the friction introduced by the \ac{DP}, which was raised by the experts.

\designprinciple{DP6}{Establish request purpose before querying}
    {For the \textit{developer} to \textit{tune model behaviour} to a \textit{user}'s needs}
    {use a \textit{role taxonomy} to establish the collaboration mode with the user before the model responds.}
    {Because users often have clear expectations about the system's role and responses, that are not always clearly communicated to the system. Forcing users to implicitly (or explicitly) communicate their expectations before the system responds, improves the alignment of expectations.}

\textbf{DP-F} is reworded completely to overcome the \textit{accessibility} issues identified during the interviews. The revised \textbf{DP7} uses plainer language omitting the concept of cognitive offloading in the mechanism and instead describing the offloading process.

\designprinciple{DP7}{Provide a shared notes space}
    {For the \textit{developer} to provide \textit{shared working surfaces} for the \textit{user}}
    {offer (task specific) \textit{inputs and outputs} separate from the main conversation.}
    {Because users will use externalisation / cognitive offloading, offering a system readable working surface reduces the cognitive effort by the users to track and communicate problem information to the system.}

Three \acp{DP} (\textbf{DP-B}, \textbf{DP-E}, \textbf{DP-F}) were reworded to improve \textit{accessibility}, without changing their mechanism, rationale, or \ac{DR}-mapping. The other three (\textbf{DP-A}, \textbf{DP-C}, \textbf{DP-D}) were revised, including an update to the \ac{DR}-mapping. Additionally, the sycophancy risk raised by the experts for \textbf{DP-A} is addressed in the revised set. Furthermore, due to unprompted expert feedback, the wording of \textbf{DR1} was also revised replacing \enquote{intelligibility} with \enquote{understandability}. The revised set of \acp{DP} is shown in Figure~\ref{fig:mr-dr-dp-mapping-revised}, which highlights rewording changes in blue and revised \ac{DP} in purple.

\begin{figure}[ht]
    \centering
    \resizebox{\textwidth}{!}{%
    \begin{tikzpicture}[
        node distance=0pt,
        mrbox/.style={neutralnode, align=left, text width=5.1cm},
        drbox/.style={neutralnode, align=left, text width=6cm},
        drboxword/.style={bluenode, align=left, text width=6cm},
        dpbox/.style={neutralnode, align=left, text width=6cm},
        dpboxword/.style={bluenode, align=left, text width=6cm},
        dpboxrev/.style={violetnode, align=left, text width=6cm},
        mr1line/.style={paletteblue, line width=1pt},
        mr2line/.style={palettered, line width=1pt},
        mr3line/.style={paletteochre, line width=1pt},
        dr1line/.style={palettegreen, line width=1pt},
        dr2line/.style={paletteochre, line width=1pt},
        dr3line/.style={paletteviolet, line width=1pt},
        dr4line/.style={paletteblue, line width=1pt},
        dr5line/.style={palettered, line width=1pt},
        partial/.style={densely dotted},
    ]
        % ---- Meta-requirements ----
        \node[mrbox] (mr1) at (0,2.3)  {\textbf{MR1:}~Support accurate mental models};
        \node[mrbox] (mr2) at (0,0)    {\textbf{MR2:}~Enable appropriate reliance};
        \node[mrbox] (mr3) at (0,-2.3) {\textbf{MR3:}~Support appropriate cognitive engagement};

        % ---- Design requirements ----
        \node[drboxword] (dr1) at (8,3.0)  {\textbf{DR1:}~Increase understandability of system behaviour};
        \node[drbox] (dr2) at (8,1.5)  {\textbf{DR2:}~Align capability expectations with actual system capability};
        \node[drbox] (dr3) at (8,0.0)  {\textbf{DR3:}~Increase verifiability of system outputs};
        \node[drbox] (dr4) at (8,-1.5) {\textbf{DR4:}~Increase deliberate engagement at consequential decision points};
        \node[drbox] (dr5) at (8,-3.0) {\textbf{DR5:}~Decrease cognitive effort of interpreting system outputs};

        % ---- Design principles (revised, seven total) ----
        \node[dpboxrev] (dp1)  at (17,4.0)  {\textbf{DP1:}~Adapt responses to user expertise};
        \node[dpboxword] (dp2)  at (17,2.6)  {\textbf{DP2:}~Disclose system limitations};
        \node[dpboxrev] (dp3) at (17,1.2)  {\textbf{DP3:}~Indicate information provenance};
        \node[dpboxrev] (dp4) at (17,-0.05)  {\textbf{DP4:}~Surface conflicting information};
        \node[dpboxrev] (dp5)  at (17,-1.4) {\textbf{DP5:}~Make system assumptions inspectable and correctable};
        \node[dpboxword] (dp6)  at (17,-2.8) {\textbf{DP6:}~Establish request purpose before querying};
        \node[dpboxword] (dp7)  at (17,-4.05) {\textbf{DP7:}~Provide a shared notes space};

        % ---- Column headings ----
        \node[font=\bfseries] at (0.0,5.0) {Meta-Requirements};
        \node[font=\bfseries] at (8.0,5.0) {Design Requirements};
        \node[font=\bfseries] at (17.0,5.0) {Design Principles};

        % ---- MR -> DR edges, unchanged from Figure~\ref{fig:mr-dr-dp-mapping} ----
        \draw[mr1line]          (mr1.east) -- (dr1.west);
        \draw[mr1line]          (mr1.east) -- (dr2.west);
        \draw[mr2line]          (mr2.east) -- (dr1.west);
        \draw[mr2line]          (mr2.east) -- (dr2.west);
        \draw[mr2line]          (mr2.east) -- (dr3.west);
        \draw[mr2line]          (mr2.east) -- (dr4.west);
        \draw[mr3line]          (mr3.east) -- (dr1.west);
        \draw[mr3line,partial]  (mr3.east) -- (dr3.west);
        \draw[mr3line]          (mr3.east) -- (dr4.west);
        \draw[mr3line]          (mr3.east) -- (dr5.west);

        % ---- DR -> DP edges, revised ----
        \draw[dr1line]          (dr1.east) -- (dp1.west);
        \draw[dr5line]          (dr5.east) -- (dp1.west);
        \draw[dr2line]          (dr2.east) -- (dp2.west);
        \draw[dr1line]          (dr1.east) -- (dp3.west);
        \draw[dr3line]          (dr3.east) -- (dp3.west);
        \draw[dr4line]          (dr4.east) -- (dp4.west);
        \draw[dr1line]          (dr1.east) -- (dp5.west);
        \draw[dr3line]          (dr3.east) -- (dp5.west);   % remapped from DR1/DR5
        \draw[dr5line]          (dr5.east) -- (dp5.west);
        \draw[dr2line]          (dr2.east) -- (dp6.west);
        \draw[dr4line]          (dr4.east) -- (dp6.west);
        \draw[dr5line, partial] (dr5.east) -- (dp6.west);
        \draw[dr5line]          (dr5.east) -- (dp7.west);
    \end{tikzpicture}%
    }
    \caption[Revised design principles]{The revised design principles mapped onto the design requirements which they fulfil. Rewording changes are highlighted in blue, revised \ac{DP} in purple.}
    \label{fig:mr-dr-dp-mapping-revised}
\end{figure}

\subsection{Uncovered Design Space} \label{ssec:uncovered-design-space}

This section discusses the design space that was not covered by the \acp{DP} derived in Section~\ref{ssec:initial-design-principles}. Firstly, potential reasons for the lack of empirically grounded \acp{DP} are discussed in Section~\ref{sssec:gap-analysis}. Secondly, a set of \acp{DP} based on the literature discussed in Chapters~\ref{sec:fundamentals} and \ref{sec:problem-awareness} is proposed to close the identified gaps in Section~\ref{sssec:theory-dp}.

\subsubsection{Gap Analysis} \label{sssec:gap-analysis}

The design theory presented in Section~\ref{ssec:design-principle-revision} leaves two gaps in the design space. Firstly, the \acp{DP} derived from the empirical data do not address explanations directly. The only related finding, is the difference in heuristic uptake of \ac{LLM} responses between the two groups. Secondly, \textbf{DR2} is only addressed one-sidedly with \textbf{DP2} prescribing the disclosure of limitations, but no \ac{DP} addresses the proactive communication of capabilities. There are several potential reasons for these gaps which are discussed below.

The first potential reason, is that the task complexity was too low to induce the need for explanation engagement. This is supported by the coding results showing that participants' effort was focused on solving the problems themselves rather than engaging with the \ac{LLM}: \textit{Problem Cognition} is the largest aggregate dimension by statement count. At the chosen difficulty level, engaging with the explanation (if provided) was optional rather than necessary, leading to the small number of explanation-engagements for group E. This hypothesis is also supported by the fact, that 32 of 36 tasks worked on by participants were solved correctly. Of the remaining four tasks two were solved incorrectly, and two were aborted and replaced with a new task. Parts of the low engagement with explanations could also be traced back to the UI design, which rendered the explanations collapsed by default. P04 noted during the debriefing that explanations should be expanded by default.

Another potential reason, is that the methodological approach could have missed silent engagement with explanations. The think-aloud protocols evaluated for this study may not have captured silent reading (and reflection) of explanations. \textcite{Ericsson2017} notes that \enquote{thoughts that pass through attention are not always verbalised}, and therefore may not be captured in the coding scheme.

The quantitative data gathered with the post-experiment questionnaires corroborate these hypotheses. Seven constructs were evaluated in total, by calculating the mean for each group and comparing them using a one-sided Mann-Whitney U test, with the null hypothesis, that group E would show higher scores for acceptance and efficacy constructs, and lower scores for workload constructs. An alpha of $\alpha=0.1$ was used throughout, to account for the small sample size and the exploratory nature of the study. The results are summarised in Table~\ref{tab:quant-results-gap}.

\begin{ctable}
    \centering
    \begin{tabular}{l c c c}
        \toprule
        \textbf{Construct} & \textbf{Mean (B)} & \textbf{Mean (E)} & \textbf{p-value} \\
        \midrule
        Perceived Usefulness $\uparrow$   & 5.233  & 5.905  & 0.095* \\
        Perceived Ease of Use $\uparrow$  & 6.300  & 5.929  & 0.745  \\
        Computer Self-Efficacy $\uparrow$ & 9.240  & 8.214  & 0.659  \\
        Mental Load $\downarrow$          & 12.800 & 10.714 & 0.282  \\
        Performance $\downarrow$          & 8.000  & 3.857  & 0.143  \\
        Effort $\downarrow$               & 10.800 & 9.429  & 0.500  \\
        Frustration $\downarrow$          & 12.200 & 2.571  & 0.005* \\
        \bottomrule
    \end{tabular}
    \caption[Quantitative results of the study]{Quantitative comparison of Group B and Group E. Arrows indicate the direction in which the explanation-augmented condition (Group E) was hypothesised to differ in the original experiment. Constructs marked with * are significant at $\alpha=0.1$.}
    \label{tab:quant-results-gap}
\end{ctable}

Only two constructs show significant differences. \textbf{Frustration} shows a large effect, indicating that the explanation-augmentation of the experiment worked in general. Frustration, however, is a workload construct and speaks to how the interaction felt, not to the user's mental model. \textbf{Perceived Usefulness}, which is an acceptance construct, shows a smaller effect and is only significant at $\alpha=0.1$. Because it reflects the user's perception of the system's utility it is affected by the user's mental model. The effect therefore provides additional support for the silent-reading hypothesis. The remaining constructs show no significant differences, but the workload constructs show a directional effect, with group E reporting lower workload than group B. The two remaining acceptance and efficacy constructs show effects against the hypothesised direction. For \textbf{Perceived Ease of Use} this could be explained by the additional interface element present only for group E. \textbf{Computer Self-Efficacy} measures a dispositional trait, which could be an indication of imperfect group equivalence, rather than a result of the experiment.

\subsubsection{Theory-derived Design Principles} \label{sssec:theory-dp}

This section proposes \ac{DP} derived from the literature reviewed in Chapters~\ref{sec:fundamentals} and \ref{sec:problem-awareness}, to close the gaps identified above. Because the \acp{DP} are developed after the expert interviews concluded they are not evaluated in the same way as \textbf{DP1}~--~\textbf{DP7} and are numbered separately to avoid ambiguity. The same modified Gregor format is used and mappings to the \acp{DR} are provided.

\textbf{DR1} calls for the system to be understandable, with \textbf{DP1}, \textbf{DP3}, and \textbf{DP5} addressing aspects of the requirement. \textbf{DP3} and \textbf{DP5} address the \textit{knowledge distribution} aspect of mental models and \textbf{DP2} and \textbf{DP-T5} seek to influence the model of \textit{global behaviour} \parencite{Gero2020}. None of the \acp{DP} address the \textit{local behaviour} aspect, which is particularly important for generative systems, where the same request can yield materially different outputs on repeated invocation \parencite{Weisz2025}. Therefore, \textbf{DP-T1} is proposed to address this gap. Providing local explanations addresses \textbf{DR1}. If the explanations are faithful they can also support \textbf{DR3}, as they allow the user to verify the system's reasoning process.

\designprinciple[palettegrey]{DP-T1}{Explain local behaviour}
    {For the \textit{developer} to provide \textit{explanations of outputs} to \textit{users}}
    {accompany outputs with an account of \textit{why this specific response} was produced.}
    {Because generative variability means a general account of system behaviour does not explain any single output, explaining local behaviour is crucial to make system behaviour understandable and verifiable.}

\ac{CoT} explanations are not guaranteed to reflect a system's underlying reasoning process i.e., be unfaithful to the systems' reasoning process \parencite{Lyu2023}. An unfaithful explanation corrupts the user's mental model and can lead to miscalibrated reliance on the system's output \parencite{Jussupow2021}. Because guaranteeing faithfulness is not a solved problem \parencite{Tutek2025}, \textbf{DP-T2} aims to mitigate the risk of unfaithful explanations. By treating explanations as challengeable evidence this \ac{DP} supports \textbf{DR1} and \textbf{DR3}, even where faithfulness cannot be guaranteed.

\designprinciple[palettegrey]{DP-T2}{Account for unfaithful explanations}
    {For the \textit{developer} to design for the \textit{possibility of unfaithful explanations} provided to \textit{users}}
    {design systems that encourage the user to \textit{treat explanations as challengeable evidence of the system's process}, rather than as a definitive account of the system's reasoning.}
    {Because an unfaithful explanation corrupts the user's mental model and can lead to miscalibrated reliance.}

The timing of explanation delivery should be oriented towards the user's curiosity and need. G3 by \textcite{Amershi2019} already calls for systems to \enquote{time services based on context}, and \textbf{DP-T3} applies this to explanations specifically. Explanations should be provided when the user encounters unexpected behaviour or their mental model is challenged, not preemptively \parencite{Miller2019,Hoffman2023}, because early or indiscriminate delivery risks anchoring the user's assessment before independent consideration \parencite{Jussupow2021}. The mechanism for this \ac{DP} is aspirational rather than directly actionable, as it requires the system to model the user's mental state in real time. \ac{MToM} seeks to address this open problem, but while theoretical frameworks exist, validated implementations are not yet available  \parencite{Wang2021, Wang2024a}. Well-timed explanations support deliberate engagement at crucial points (\textbf{DR4}) and indirectly assist forming an understanding of system behaviour (\textbf{DR1}).

\designprinciple[palettegrey]{DP-T3}{Time explanation delivery}
    {For the \textit{developer} to \textit{time the delivery of explanations} for \textit{users}}
    {by providing explanations \textit{when the user's curiosity is raised} or \textit{their mental model is challenged}.}
    {Because early delivery risks anchoring the user's assessment before an independent one can form. Furthermore, indiscriminate delivery risks overwhelming the user with information.}

An explanation can supply the material needed by users for an unwarranted dismissal of system output driving the \enquote{explaining away} pattern described by \textcite{Jussupow2021}. Because the user is already engaged with the system's output and explanation, merely inducing engagement does not address this risk. Instead, this risk should be addressed by introducing additional friction at the time of dismissal \parencite[cf.][]{Buinca2021}. A potential implementation of this principle is to require the user to articulate a reason rooted in the facts outside the explanation before overriding or dismissing an output. This mechanism is similar to the friction \textbf{DP6} introduces at the point of query, but is applied at the point of dismissal. This \ac{DP} supports deliberate engagement (\textbf{DR4}).

\designprinciple[palettegrey]{DP-T4}{Constrain post-hoc dismissal}
    {For the \textit{developer} to make \textit{unwarranted dismissal of a system's output} harder for \textit{users}}
    {require the user to \textit{articulate a fact-based reason for dismissal} before overriding or dismissing an output.}
    {Because explanations can provide a rationale for \enquote{explaining away} the systems outputs and users should dismiss deliberately.}

The second gap identified in Section~\ref{sssec:gap-analysis} was the lack of \acp{DP} addressing the communication of system capabilities. \textbf{DP-T5} addresses this gap building on G1 and G2 by \textcite{Amershi2019}, which call to disclose the system's capabilities. A key reason, why this was not addressed in the empirical data, is that the communication of capabilities should happen \enquote{initially} \parencite{Amershi2019}. The experiment design did include an initial briefing, which was not captured in the transcripts, and therefore not coded. The \ac{DP} complements \textbf{DP2} by proactively communicating capabilities. The \ac{DP} addresses \textbf{DR2} by ensuring that users are aware of the system's capabilities.

\designprinciple[palettegrey]{DP-T5}{Proactively communicate system capabilities}
    {For the \textit{developer} to \textit{communicate capabilities} to the \textit{users}}
    {\textit{brief users on the system's capabilities} before the first interaction.}
    {Because communicating capabilities accurately assists in building the users initial mental model of the system.}

Beyond the five principles above, it is important to note that \textbf{DP1} itself can be extended to cover explanations specifically. The same grounding already cited for \textbf{DP1} shows that a given explanation can achieve different levels of causability for different users \parencite{Shin2021}, and tailoring explanations to the user aligns with the audience-specific definition of \ac{XAI} adopted in Section~\ref{sssec:terminology} \parencite{Arrieta2020}.

Figure~\ref{fig:mr-dr-dp-mapping-final} consolidates the complete design theory, combining the seven empirically grounded and expert-evaluated principles from Section~\ref{sssec:revised-dp} with the five theory-derived principles proposed above.

\begin{figure}[ht]
    \centering
    \resizebox{\textwidth}{!}{%
    \begin{tikzpicture}[
        node distance=0pt,
        mrbox/.style={neutralnode, align=left, text width=5.1cm},
        drbox/.style={neutralnode, align=left, text width=6cm},
        dpbox/.style={neutralnode, align=left, text width=6cm},
        dptbox/.style={bluenode, align=left, text width=6cm},
        mr1line/.style={paletteblue, line width=1pt},
        mr2line/.style={palettered, line width=1pt},
        mr3line/.style={paletteochre, line width=1pt},
        dr1line/.style={palettegreen, line width=1pt},
        dr2line/.style={paletteochre, line width=1pt},
        dr3line/.style={paletteviolet, line width=1pt},
        dr4line/.style={paletteblue, line width=1pt},
        dr5line/.style={palettered, line width=1pt},
        partial/.style={densely dotted},
    ]
        % ---- Meta-requirements ----
        \node[mrbox] (mr1) at (0,2.3)  {\textbf{MR1:}~Support accurate mental models};
        \node[mrbox] (mr2) at (0,0)    {\textbf{MR2:}~Enable appropriate reliance};
        \node[mrbox] (mr3) at (0,-2.3) {\textbf{MR3:}~Support appropriate cognitive engagement};

        % ---- Design requirements ----
        \node[drbox] (dr1) at (8,4.0)  {\textbf{DR1:}~Increase understandability of system behaviour};
        \node[drbox] (dr2) at (8,2.0)  {\textbf{DR2:}~Align capability expectations with actual system capability};
        \node[drbox] (dr3) at (8,0.0)  {\textbf{DR3:}~Increase verifiability of system outputs};
        \node[drbox] (dr4) at (8,-2.0) {\textbf{DR4:}~Increase deliberate engagement at consequential decision points};
        \node[drbox] (dr5) at (8,-4.0) {\textbf{DR5:}~Decrease cognitive effort of interpreting system outputs};

        % ---- Design principles ----
        \node[dptbox] (dpt1) at (17,6.40)  {\textbf{DP-T1:}~Explain local behaviour};
        \node[dpbox]  (dp2)  at (17,5.4)  {\textbf{DP2:}~Disclose system limitations};
        \node[dptbox] (dpt5) at (17,4.20)  {\textbf{DP-T5:}~Proactively communicate system capabilities};
        \node[dpbox]  (dp3)  at (17,2.85)  {\textbf{DP3:}~Indicate information provenance};
        \node[dptbox] (dpt2) at (17,1.50)  {\textbf{DP-T2:}~Account for unfaithful explanations};
        \node[dptbox] (dpt3) at (17,0.30)  {\textbf{DP-T3:}~Time explanation delivery};
        \node[dpbox]  (dp5)  at (17,-0.90) {\textbf{DP5:}~Make system assumptions inspectable and correctable};
        \node[dptbox] (dpt4) at (17,-2.10) {\textbf{DP-T4:}~Constrain post-hoc dismissal};
        \node[dpbox]  (dp1)  at (17,-3.25) {\textbf{DP1:}~Adapt responses to user expertise};
        \node[dpbox]  (dp4)  at (17,-4.40) {\textbf{DP4:}~Surface conflicting information};
        \node[dpbox]  (dp6)  at (17,-5.60) {\textbf{DP6:}~Establish request purpose before querying};
        \node[dpbox]  (dp7)  at (17,-6.80) {\textbf{DP7:}~Provide a shared notes space};

        % ---- Column headings ----
        \node[font=\bfseries] at (0.0,7.4)  {Meta-Requirements};
        \node[font=\bfseries] at (8.0,7.4)  {Design Requirements};
        \node[font=\bfseries] at (17.0,7.4) {Design Principles};

        % ---- MR -> DR edges ----
        \draw[mr1line]          (mr1.east) -- (dr1.west);
        \draw[mr1line]          (mr1.east) -- (dr2.west);
        \draw[mr2line]          (mr2.east) -- (dr1.west);
        \draw[mr2line]          (mr2.east) -- (dr2.west);
        \draw[mr2line]          (mr2.east) -- (dr3.west);
        \draw[mr2line]          (mr2.east) -- (dr4.west);
        \draw[mr3line]          (mr3.east) -- (dr1.west);
        \draw[mr3line,partial]  (mr3.east) -- (dr3.west);
        \draw[mr3line]          (mr3.east) -- (dr4.west);
        \draw[mr3line]          (mr3.east) -- (dr5.west);

        % ---- DR -> DP edges ----
        \draw[dr1line]          (dr1.east) -- (dp1.west);
        \draw[dr5line]          (dr5.east) -- (dp1.west);
        \draw[dr2line]          (dr2.east) -- (dp2.west);
        \draw[dr1line]          (dr1.east) -- (dp3.west);
        \draw[dr3line]          (dr3.east) -- (dp3.west);
        \draw[dr4line]          (dr4.east) -- (dp4.west);
        \draw[dr1line]          (dr1.east) -- (dp5.west);
        \draw[dr3line]          (dr3.east) -- (dp5.west);
        \draw[dr5line]          (dr5.east) -- (dp5.west);
        \draw[dr2line]          (dr2.east) -- (dp6.west);
        \draw[dr4line]          (dr4.east) -- (dp6.west);
        \draw[dr5line, partial] (dr5.east) -- (dp6.west);
        \draw[dr5line]          (dr5.east) -- (dp7.west);

        \draw[dr1line]          (dr1.east) -- (dpt1.west);
        \draw[dr1line]          (dr1.east) -- (dpt2.west);
        \draw[dr3line, partial] (dr3.east) -- (dpt1.west);
        \draw[dr3line]          (dr3.east) -- (dpt2.west);
        \draw[dr4line]          (dr4.east) -- (dpt3.west);
        \draw[dr1line, partial] (dr1.east) -- (dpt3.west);
        \draw[dr4line]          (dr4.east) -- (dpt4.west);
        \draw[dr2line]          (dr2.east) -- (dpt5.west);
    \end{tikzpicture}%
    }
    \caption[Final design theory, including theory-derived design principles]{The complete design theory. \acp{DP} are organised to reduce the number of edge crossings. Theory-derived principles are highlighted in blue.}
    \label{fig:mr-dr-dp-mapping-final}
\end{figure}
